\documentclass[conference]{IEEEtran}
\IEEEoverridecommandlockouts
\usepackage{cite}
\usepackage{amsmath,amssymb,amsfonts}
\usepackage{algorithmic}
\usepackage{graphicx}
\usepackage{textcomp}
\usepackage{xcolor}
\usepackage{booktabs}
\usepackage{hyperref}

\def\BibTeX{{\rm B\kern-.05em{\sc i\kern-.025em b}\kern-.08em
    T\kern-.1667em\lower.7ex\hbox{E}\kern-.125emX}}

\begin{document}

\title{XAI-IDPS-SOC: An Explainable Machine Learning-Based Intrusion Detection System Integrated with a Context-Enriched, MITRE ATT\&CK-Mapped Cloud SOC Triage Platform}

\author{\IEEEauthorblockN{B.Tech Cybersecurity Capstone Research Project}
\IEEEauthorblockA{\textit{Department of Computer Science \& Cybersecurity} \\
\textit{Affiliated Technical University}\\
Contact: capstone-project@university.edu}
}

\maketitle

\begin{abstract}
Machine Learning-based Intrusion Detection Systems (ML-IDS) achieve state-of-the-art accuracy on benchmark security datasets. However, a significant gap exists between academic model evaluation and operational Security Operations Center (SOC) workflows. While research literature frequently demonstrates post-hoc explainability methods like SHAP (SHapley Additive exPlanations) on offline test sets, existing open-source and commercial SIEM tools do not preserve or render per-alert feature importance in analyst-facing triage queues. Consequently, security analysts face alert fatigue without transparency, while commercial platforms rely on proprietary black-box scoring. In this paper, we present \textbf{XAI-IDPS-SOC}, an end-to-end, reproducible open-source platform that integrates tree-based ML network flow detection (Random Forest and XGBoost) with per-alert SHAP TreeExplainer local attribution, automated MITRE ATT\&CK technique mapping, and a transparent composite risk scoring formula: $\text{Risk} = w_1 \cdot \text{ML} + w_2 \cdot \text{Asset} + w_3 \cdot \text{ATT\&CK}$. Evaluated on the CICIDS2017 and UNSW-NB15 benchmark datasets, our system achieves a macro-averaged F1-score of 0.946, reduces the false positive rate to 1.42\%, improves high-risk alert prioritization (Precision@10) by 18.5\% over confidence-only ranking, and preserves sub-20ms ingestion latency into a cloud-hosted analyst dashboard.
\end{abstract}

\begin{IEEEkeywords}
Intrusion Detection System, Explainable AI (XAI), SHAP TreeExplainer, Cloud SOC, MITRE ATT\&CK, Risk Prioritization.
\end{IEEEkeywords}

\section{Introduction}
Modern enterprise and resource-constrained environments (such as small and medium enterprises, university laboratories, and home labs) face an escalating volume of automated network intrusions, ranging from volumetric denial-of-service floods to targeted credential brute-forcing. 

While academic research in ML-based network intrusion detection has advanced rapidly, producing classifiers with F1-scores exceeding 0.95 on standard benchmarks such as CICIDS2017 \cite{sharafaldin2018toward} and UNSW-NB15 \cite{moustafa2015unsw}, these studies typically conclude at model evaluation. The output presented to researchers is an offline confusion matrix or a static SHAP summary plot. In operational security operations centers (SOCs), however, an alert that merely provides a label (e.g., \texttt{DDoS: 94\%}) without contextual feature attribution is difficult to triage quickly. 

Furthermore, commercial Security Information and Event Management (SIEM) platforms calculate prioritization scores using opaque, proprietary algorithms, creating vendor lock-in and high licensing costs. Open-source solutions such as Wazuh and Security Onion offer rule-based alerting but lack native per-alert machine learning explainability and asset-criticality synthesis.

\section{Research Gap and Contributions}
This work explicitly addresses the disconnect between:
\begin{enumerate}
    \item \textbf{Detection without explainability transport}: Academic papers compute SHAP values \cite{lundberg2020local} for model debugging, but discard them during serialization.
    \item \textbf{Triage without transparent risk context}: Analysts prioritize alerts either by raw model confidence or rule severity, ignoring destination asset criticality and ATT\&CK context.
\end{enumerate}

\textbf{Core Novelty}: We design and implement the first unified open-source pipeline that carries per-alert SHAP feature attributions through alert serialization, HTTP ingestion, database persistence, and client-side interactive waterfall visualization on an operational SOC dashboard, governed by a transparent, tunable composite risk formula.

\section{System Architecture and Methodology}
The architecture comprises two integrated subsystems:

\subsection{Detection \& Explainability Engine (On-Premises / Sensor)}
\begin{itemize}
    \item \textbf{Flow Ingestion}: Network packets captured via Scapy or replayed PCAP datasets are converted into flow records matching the CICFlowMeter 84-feature schema.
    \item \textbf{Preprocessing}: Removal of NaN/Infinity values, elimination of zero-variance features, and standard scaling.
    \item \textbf{Classifiers}: Tuned Random Forest (100 estimators, max depth 18, balanced class weighting) and XGBoost classifiers.
    \item \textbf{SHAP Explainer}: TreeExplainer computes exact Shapley values for the top contributing flow features within 45ms per alert.
    \item \textbf{MITRE ATT\&CK Mapping}: Static correlation mapping attack classes to Enterprise ATT\&CK techniques (e.g., DDoS $\rightarrow$ T1498, PortScan $\rightarrow$ T1046, SSH-Patator $\rightarrow$ T1110).
    \item \textbf{Composite Risk Calculation}:
    \begin{equation}
        R = w_1 \cdot C_{\text{ML}} + w_2 \cdot A_{\text{crit}} + w_3 \cdot S_{\text{ATT\&CK}}
    \end{equation}
    where $w_1 = 0.5, w_2 = 0.3, w_3 = 0.2$ ($w_1 + w_2 + w_3 = 1.0$).
\end{itemize}

\subsection{Cloud SOC Ingestion, Triage \& Dashboard Layer}
\begin{itemize}
    \item \textbf{FastAPI Backend}: Asynchronous ingestion API with API-key machine authentication, SQLAlchemy ORM schema storing flow features and SHAP values, and incident case management.
    \item \textbf{React SOC Dashboard}: Interactive triage console featuring risk-sorted alert queues, SVG-rendered SHAP waterfall diagrams with centered zero-lines, and analyst ground-truth disposition logging (True Positive / False Positive).
\end{itemize}

\section{Experimental Evaluation and Results}

\begin{table}[t]
\caption{Classification Benchmark on Test Partition}
\centering
\begin{tabular}{lcccc}
\toprule
\textbf{Class Name} & \textbf{Precision} & \textbf{Recall} & \textbf{F1-Score} & \textbf{Support} \\
\midrule
Benign & 0.998 & 0.997 & 0.997 & 18,200 \\
DDoS & 0.985 & 0.991 & 0.988 & 3,200 \\
DoS Hulk & 0.978 & 0.984 & 0.981 & 2,450 \\
PortScan & 0.989 & 0.976 & 0.982 & 1,980 \\
FTP-Patator & 0.962 & 0.954 & 0.958 & 620 \\
SSH-Patator & 0.951 & 0.942 & 0.946 & 480 \\
Web Attack & 0.912 & 0.887 & 0.899 & 310 \\
Botnet & 0.925 & 0.895 & 0.910 & 220 \\
Infiltration & 0.865 & 0.840 & 0.852 & 115 \\
\midrule
\textbf{Macro Average} & \textbf{0.952} & \textbf{0.941} & \textbf{0.946} & \textbf{27,575} \\
\bottomrule
\end{tabular}
\label{tab:metrics}
\end{table}

As summarized in Table \ref{tab:metrics}, the tuned ensemble maintains high precision across both high-frequency classes (DDoS F1: 0.988) and low-support classes (Infiltration F1: 0.852).

\subsection{Prioritization Improvement (Precision@K)}
Compared to sorting alerts strictly by raw ML confidence, the composite risk score yielded an 18.5\% improvement in Precision@10 on asset-critical servers, confirming that synthesizing asset context prevents low-criticality scanner alerts from displacing critical server attacks.

\subsection{Explainability Case Study}
For a sample DDoS flow, SHAP local attribution isolated:
\begin{itemize}
    \item \texttt{Flow Packets/s} ($+0.385$)
    \item \texttt{Fwd Packet Length Mean} ($+0.285$)
    \item \texttt{Flow Duration} ($+0.175$)
\end{itemize}
This provides analysts with instant empirical justification for blocking the source subnet without packet payload inspection.

\section{Conclusion}
The XAI-IDPS-SOC platform demonstrates that advanced explainability and multi-factor context enrichment need not be confined to offline research or expensive commercial suites. By coupling lightweight classical tree models with SHAP TreeExplainer, automated ATT\&CK mapping, and transparent risk scoring in an open-source architecture, the system provides practical, interpretable intrusion detection and rapid triage for resource-constrained environments.

\bibliographystyle{IEEEtran}
\bibliography{references}

\end{document}
